\documentclass[12pt]{article}
\usepackage{amsmath,amssymb,amsfonts}
\usepackage[margin=1in]{geometry}

\begin{document}

\begin{center}
{\Large\bfseries Chapter 10, Problem 1}
\end{center}

\noindent\textbf{Problem.} Consider the plane of Fig.~10.3, where the objects at $A$, $B$, and $C$ are fixed to the circular loop. This figure is clearly invariant under rotation through $120^\circ$ and $240^\circ$ and under reflection through the lines $a$, $b$, and $c$. Show that these five operations plus the identity form a group. Write out the group multiplication table. What are the classes into which the various group elements fall?

\bigskip
\noindent\textbf{Solution.}

The six symmetry operations of the equilateral triangle are:
\begin{itemize}
\item $e$: the identity
\item $R_{120}$: rotation by $120^\circ$ (counterclockwise)
\item $R_{240}$: rotation by $240^\circ$ (counterclockwise)
\item $\sigma_a$: reflection through line $a$ (through vertex $A$)
\item $\sigma_b$: reflection through line $b$ (through vertex $B$)
\item $\sigma_c$: reflection through line $c$ (through vertex $C$)
\end{itemize}

This is the group $S_3$ (the symmetric group on three objects), also called the dihedral group $D_3$.

\medskip
\noindent\textbf{Verification of group axioms:}

\noindent (1) \textit{Closure:} The composition of any two symmetry operations is again a symmetry operation, which can be verified by computing all 36 products.

\noindent (2) \textit{Associativity:} Composition of functions is always associative.

\noindent (3) \textit{Identity:} The identity operation $e$ satisfies $e \cdot g = g \cdot e = g$ for all $g$.

\noindent (4) \textit{Inverses:} $e^{-1} = e$, $R_{120}^{-1} = R_{240}$, $R_{240}^{-1} = R_{120}$, $\sigma_a^{-1} = \sigma_a$, $\sigma_b^{-1} = \sigma_b$, $\sigma_c^{-1} = \sigma_c$ (reflections are their own inverses).

\medskip
\noindent\textbf{Multiplication Table:}

We label the elements as $e, R, R^2, \sigma_a, \sigma_b, \sigma_c$ where $R = R_{120}$ and $R^2 = R_{240}$.

\[
\begin{array}{c|cccccc}
 & e & R & R^2 & \sigma_a & \sigma_b & \sigma_c \\
\hline
e & e & R & R^2 & \sigma_a & \sigma_b & \sigma_c \\
R & R & R^2 & e & \sigma_c & \sigma_a & \sigma_b \\
R^2 & R^2 & e & R & \sigma_b & \sigma_c & \sigma_a \\
\sigma_a & \sigma_a & \sigma_b & \sigma_c & e & R & R^2 \\
\sigma_b & \sigma_b & \sigma_c & \sigma_a & R^2 & e & R \\
\sigma_c & \sigma_c & \sigma_a & \sigma_b & R & R^2 & e
\end{array}
\]

\medskip
\noindent\textbf{Classes (conjugacy classes):}

Two elements $g_1$ and $g_2$ are in the same class if there exists a group element $h$ such that $g_2 = h g_1 h^{-1}$. Computing:

\begin{enumerate}
\item $\{e\}$: The identity is always in a class by itself.

\item $\{R, R^2\}$: We have $\sigma_a R \sigma_a^{-1} = \sigma_a R \sigma_a = R^2$, so $R$ and $R^2$ are conjugate.

\item $\{\sigma_a, \sigma_b, \sigma_c\}$: We have $R \sigma_a R^{-1} = R \sigma_a R^2 = \sigma_b$ and $R^2 \sigma_a R = \sigma_c$, so all three reflections are conjugate.
\end{enumerate}

Thus the group has three conjugacy classes: $\{e\}$, $\{R, R^2\}$, and $\{\sigma_a, \sigma_b, \sigma_c\}$.

\end{document}
